We propose \textbf{Actionable Specificity (AS)}, a framework decomposing verification signal informativeness into localization, state exposure, and causal context to explain why different feedback types yield different LLM code repair success rates. Testing AS extractability on HumanEval and MBPP (664 problems, 600 signals across 6 conditions), we find regex-based extraction achieves only \textbf{49.5\% coverage}, with state information at \textbf{2.7\%}. \textbf{This is a limitation of our operationalization}: assertion errors dominate these benchmarks (\textbf{53.4\%} of failures) and produce minimal traceback output lacking the file:line format our extraction patterns require. Prior work suggests assertion errors have among the lowest repair success rates, consistent with limited feedback availability. While extraction failed to meet our 95\% target, AS ordering (C1$\geq$C2$\geq$C3$\geq$C4) holds on the extractable subset, and the failure itself reveals a \textbf{boundary condition}: feedback informativeness research must verify extraction feasibility by error type before operationalizing measurement.
